
\subsection{Fantom}

Fantom uses a new consensus mechanism that they call Lachesis.
The core elements of Lachesis are: asynchronous BFT, a block-DAG, and some form of PoS.

The Fantom Foundation website says that Fantom (via Lachesis, their consensus mechanism) solves the trilemma:

% \footnote{
%     See: \url{\linkFantomIntro}, \url{\linkFantomIndexingSoln}, \url{\linkFantomLachesis}, \url{\linkFantomFaq}
% }

% \bquote{
%     Current solutions make trade-offs between three components: scalability, security and decentralization. This is known as the blockchain trilemma.

%     [...]

%     Fantom tackles the problem at the core: its high-speed consensus mechanism, Lachesis, allows digital assets to operate at unprecedented speed and delivers dramatic improvements over the current systems.
% }[\url{\linkFantomIntro}]


\bquote{
    It solves the blockchain Scalability Trilemma, [...]
}[\citeFantomLachesis]

Fantom also uses quotes like this in promotional material:

\bquote{
    \emph{``Fantom is an exceptional platform that is cutting no corners in tackling the blockchain trilemma. {[...]}''}
    -- Ganesh, Covalent CEO
}[{\href{\linkFantomIndexingSoln}{Fantom press release, 2021-03-23}}]

\subsubsection{No Citations, Reasoning, or Explanation}

Pages on Fantom's website provide no citation or reference that leads to a proper argument that Lachesis is a solution to the trilemma.
Where there is some explanation, it is superficial and insufficient.\footnote{
    Example: \citeFantomLachesis.
}

Neither the Fantom whitepaper nor their paper on Lachesis even mention ``trilemma''.\footnote{
    See: \url{\linkFantomWP} and \url{\linkFantomWPLachesis}.
}

Why isn't the claim ever accompanied by a citation?



\subsubsection{Lachesis is $O(c)$}

Fantom claim that the \emph{asynchronous} nature of Lachesis is the reason that it is scalable:

\bquote{
    Finally, aBFT networks allow for greater scalability and decentralization since there isn't excessive communication to limit the number of participating nodes.
}[{\citeFantomLachesis}]

While this does argue that aBFT is superior to pBFT, it isn't relevant to the trilemma --- communication overhead isn't the issue.

Rather, Fantom must show that Lachesis can process $> O(c)$ transactions, \emph{at a minimum}, whilst not suffering network degradation.
However, regardless of the order of blocks processed, $\nicefrac{2}{3}$ of validating nodes must agree before a block is finalized.\footnote{
    ``Transactions are finalized as soon as the nodes agree on the decision.'' --- \citeFantomLachesis.
}
Since validators are $O(c)$ constrained, Lachesis is $O(c)$ scalable.


\subsubsection{Shifting Goal Posts}

From Fantom's FAQ:

\bquote{
    The blockchain trilemma is the known trade-off in distributed ledger technologies, which have to balance between speed, security, and decentralization. \\
    {[...]}
}[{\href{\linkFantomFaq}{Answer to \emph{What problem does Fantom solve? What is the blockchain trilemma?}}}]

This is incorrect and does not represent the trilemma.

Particularly, the trilemma says nothing about \emph{speed}, rather, it's about \emph{capacity}.\footnote{
    This mistake is not made on another page of theirs: \citeFantomLachesis.
}

There are at least three qualitatively different meanings of \emph{speed} that apply to blockchains:

\begin{itemize}
    \item When a transaction is made, how long until the next block is produced (latency).
    \item When the network is congested, how long a transaction takes to be included in a block (latency under congestion).
    \item After a transaction is first included in a block, how long until that transaction is considered ``confirmed'' or ``finalized'' (confirmation time).
\end{itemize}

However, the trilemma is (canonically) about confirming $O(n) > O(c)$ transactions --- i.e., \emph{throughput} or \emph{capacity}.
None of the meanings of speed are directly relevant to solving the trilemma.
Being ``fast'' might make for a better user experience, but it's not relevant here.

% Since a solution to the trilemma should never experience global congestion, we don't care about \emph{latency under congestion}.

Misstating the trilemma in this way is changing the goal posts, which is (at best) misleading.

Moreover, transaction latency has never been the primary issue in PoW blockchains.\footnote{
    In 2014 I briefly worked on an ``instant confirmation'' prototype, for example.
    See: \linkQantaBitcointalk.
}
The primary issue has always been: \emph{how do you build a cohesive blockchain network such that $O(n)$ people can use it simultaneously?}
